
\section{Explicit sample budget and dependence on circuit width}
\label{sec:explicit-budget}

The following formulas are a sufficient bound, with no implicit dependence
on learner width or Gaussian scale. Given the target-cylinder cost $D_*$
from~\eqref{eq:standalone-cylinder-mass}, put
\begin{align*}
 \beta&=\varepsilon/4,&
 B_v&=\max\{8,8[D_*+\log(1/\beta)]\},&
 K&=D_*+1,\\
 r&=\min\{1/6,c/\sqrt{1+(L-1)B_v}\},&
 \alpha&=\varepsilon/12.
\end{align*}
Use these $K,r,\alpha$ in the compressor's parameter choices and set
$M=N_0$, $B=\max(B_{\rm small},B_{\rm large},1)$,
\[
 P=M(d+1)+(L-2)M^2+M,\qquad
 J=P\log\left[3+\frac{32B(d+L+1)(2M(B+1))^L}{r}\right].
\]
Then the explicit sufficient sample size is
\begin{equation}\label{eq:explicit-samples}
 n\ge\frac{16}{\varepsilon}\left[J+\log\frac1\delta\right].
\end{equation}
The statistical proof above is unchanged; only the value of $J$ improves.

\paragraph{Circuit width.}
For fixed $d,L,c,\varepsilon,\delta$, the target construction gives
$D_*=O(C\log^2(2+C))$. Write $F=C\log^2(2+C)$. The budgets satisfy
\[
 K,r^{-2}=O(F),\quad q,V=O(F^2),\quad m,M=O(F^3),
 \quad B=O(F^3),\quad P=O(F^6).
\]
The logarithmic grid cost is $O(\log(2+C))$ per parameter. Consequently
\[
 n=O\bigl(C^6\log^{13}(2+C)\bigr).
\]
This holds for the actual finite mean-field prior as well as every stronger
scale, uniformly over all learner widths $N\ge C$. It is an upper bound;
the argument does not show that the sixth power is necessary.

\paragraph{A conservative universal polynomial.}
Set
\[
 S=2+d+L+C+c+c^{-1}+\varepsilon^{-1}+\log(1/\delta).
\]
For universal constants, each quantity below is $O(S^p)$ with the displayed
exponent. Logarithms are absorbed only where indicated.
\begin{center}
\begin{tabular}{lr}
\toprule
Quantity & $p$\\\midrule
$D_*,K,B_v$ & 3\\
$r^{-2}$ & 6\\
$e^{-1},\eta^{-1}$ & 8, 2\\
$q,V$ & 13\\
$D_0$ & 14\\
$m,M$ & 21\\
$P$ & 43\\
$B_{\rm small}^2$, $B$ & 44, 22\\
\bottomrule
\end{tabular}
\end{center}
For example the reconstruction term $(V+q+1)/e$ is $O(S^{21})$,
the Gram term is $O(S^{14})$, and the readout term is $O(S^{12})$.
The grid uses $O(S\log S)$ bits per coordinate. Thus
$J=O(S^{44}\log S)$ and~\eqref{eq:explicit-samples} is
$O(S^{45}\log S)$, bounded by $O(S^{46})$.
This degree is a conservative multivariate envelope, not a claim of a
practical numerical guarantee. The headline theorem and degree-46 envelope now have Lean proofs, as does
the improved compressor in the form the main theorem uses (for tanh, and
without the robustness assumption only at large width). The particular numerical allocation in this section, including
its leading factor $16/\varepsilon$, is outside that certification.

\paragraph{Remaining quadratic question.}
The original prior restriction is preserved throughout this paper. It
remains open here whether its sample bound can be quadratic in target
circuit width. The current method retains $q=O(F^2)$ active directions and
samples $m=O(F^3)$ source neurons; coding the dense compressed network then
costs $O(F^6)$ parameters. A quadratic conclusion would require a further
improvement, rather than a reinterpretation of this bound.
